\section{Introduction}
\label{sec:introduction}

Semiconductor quantum dots provide discrete, atom-like carrier spectra while
retaining substantial control through composition, confinement, and external
geometry. Their neutral exciton ($X$) and biexciton ($XX$) states form the
elementary ladder for correlated photon emission. In particular, the
biexciton--exciton cascade has enabled regulated photon-pair emission
\cite{benson2000,moreau2001} and polarization-entangled photon generation
\cite{akopian2006}. The biexciton binding energy determines the separation of
the two cascade transitions and therefore affects spectral addressability,
resonant excitation, and regimes in which the emitted photons approach equal
energies.

The energy of a confined biexciton is not fixed by a single Coulomb scale.
Electron--hole attraction competes with electron--electron and hole--hole
repulsion, while the carrier masses and the confinement geometry determine how
strongly the four particles can rearrange. Consequently, biexcitons may be
bound, antibound, or close to zero binding even within the same material
system. Experimental studies have likewise shown that biexciton binding can
change strongly with confinement and dot structure
\cite{bayer1998,juska2019}. This sensitivity makes the biexciton a useful
probe of few-particle correlations, but it also makes its numerical evaluation
demanding: the binding energy is a small difference between two much larger
correlated energies.

Strongly oblate quantum dots are especially suitable for separating the roles
of axial and lateral confinement. Their two semiaxes provide independent
geometrical controls, while the large separation between axial and in-plane
energy scales permits an adiabatic treatment of the single-particle states.
Variational calculations for strongly oblate ellipsoidal dots have demonstrated
geometry-dependent energies, binding energies, recombination energies, and
radiative lifetimes \cite{hayrapetyan2019}, and subsequent work extended the
same framework to nonlinear optical response \cite{bleyan2022}. The
semi-ellipsoidal boundary considered here is relevant to flattened or
dome-like dots and has previously been used for single-particle and interband
absorption calculations \cite{hayrapetyan2012}. Unlike a full ellipsoid, the
dot occupies only one side of the basal plane, which changes the axial
boundary condition while retaining axial rotational symmetry.

The purpose of this work is to determine how the lateral semiaxis $a$ and
axial semiaxis $c$ control the corrected energy, binding, and radiative
lifetime of the biexciton in a strongly oblate
semi-ellipsoidal GaAs dot. We restrict the study to the single-particle ground
configuration and use correlated variational states for both $X$ and $XX$.
The high-dimensional integrals are evaluated after a state-matched importance
transformation. For the final biexciton calculation we retain the full
symmetrized correlation expression and remove only the redundant global
rotation, thereby avoiding the accumulation of independent quadrature errors
among analytically equivalent reduced terms. The variational parameters are
then refined by a safeguarded coordinate pattern search that verifies accepted
moves at two integration budgets.

Seven geometries form two intersecting scans: $c=0.5,1,1.5,2$ at fixed
$a=5$, and $a=3,5,7,10$ at fixed $c=1$, with all lengths measured in the
effective electron Bohr radius. From the optimized states we obtain $\EX$,
$\EXX$, $\Ebind=2\EX-\EXX$, the exciton coincidence overlap, and the
radiative lifetimes $\tauX$ and $\tauXX$. Particular attention is paid to the
numerical resolution of the small binding energy. The reported error bars are
rough propagations of the adaptive integrator's internal estimates rather than
formal statistical confidence intervals.
